% method.tex
% Method: problem setup, variance-corrected SIM composition, beta-clipping floor, R^2 branch, algorithm, properties.

We work throughout in continuously compounded annualized growth rates,
$g(t) = [\log p(t) - \log p(t-1)] / \Delta t$, where $p(t)$ is the
closing price on trading day $t$ and $\Delta t$ is the fraction of a
trading year spanned by a single trading day (we use
$\Delta t = 1/252$ in the experiments below, so the per-day log return
is $g(t)\,\Delta t$). Working in annualized growth rates rather than
raw log returns keeps variances and covariances in the units customary
in risk reporting. Growth rates compose linearly under aggregation in time and are the natural
unit for variance and covariance arithmetic, which the construction
below relies on heavily. All variances and covariances refer to
sample moments computed over the relevant time window unless stated
otherwise. The method composes a synthetic asset growth-rate path of
the form:
\begin{equation}
  g_i(t) \;=\; \alpha_i \;+\; \beta_i\, g_m(t) \;+\; \eps_i(t),
  \qquad t = 1,\dots,T,
  \label{eq:sim-composition}
\end{equation}
from an independent generator draw $\teps_i \in \R^T$ and a shared
market path $g_m \in \R^T$ over $T$ trading days by rescaling the
generator draw as $\eps_i = s_i\,\teps_i$ with a scalar
$s_i \in (0,1]$ chosen to preserve one of two per-asset marginal
targets while recovering the calibrated single-index model (SIM)
coefficients $(\alpha_i, \beta_i)$ in the large-$T$ limit.

\subsection{Problem setup}
\label{sec:method-setup}

We are given a single market growth-rate path $g_m \in \R^T$ over $T$
trading days, with sample mean $\mu_m$ and sample variance
$\sigma_m^2 = \Var[g_m]$. In the experiments of
Section~\ref{sec:results} the market path is the SPY total-return
growth-rate series; the construction is agnostic to whether $g_m$ is a
single-name proxy, a synthetic index, or a model-generated path.

For each of $N$ assets indexed by $i = 1,\dots,N$, we are given a draw
$\teps_i \in \R^T$ from a generator that produces growth rates with
the target marginal distribution: heavy tails, regime structure,
jumps, or whatever the application requires. Write
$\sigma_{\gen,i}^2 = \Var[\teps_i]$ for the sample variance of the
draw. The construction is agnostic to how the generator is fit; in our
experiments we use the per-asset hidden Markov model with Student-$t$
emissions and explicit jump components from
\texttt{JumpHMM.jl}\footnote{\url{https://github.com/varnerlab/JumpHMM.jl}},
fit to the asset's own full growth-rate history (not to residuals
after market detrending). Because the generator's variance reflects
the full return rather than the idiosyncratic part, the rescaling
$\eps_i = s_i\,\teps_i$ with
$s_i^2 = 1 - \beta_i^2 \sigma_m^2 / \sigma_{\gen,i}^2$ is exactly what
prevents the market term in the SIM composition from double-counting
variance that is already carried by $\teps_i$. The generator draws
are demeaned, $\E[\teps_i] = 0$, so that the constant term in the SIM
is carried entirely by $\alpha_i$.
We assume $\Cov[\teps_i, g_m] \approx 0$ for each $i$, which holds
when the generator and market path are independent draws (the typical
use case), and holds after a copula rank reorder when the copula is
sampled independently of $g_m$. If the generator already carries
non-trivial market correlation, the residual covariance can be
subtracted from $\beta_i$ before applying the correction; we return to
this in Section~\ref{sec:discussion}.

For each asset we are given a calibrated triple
$(\alpha_i, \beta_i, R^2_{i,\real})$, estimated by ordinary least
squares on the real return history of asset $i$ against the real
market path. These are the targets we want the synthetic construction
to recover. The intercept $\alpha_i$ and slope $\beta_i$ are the
standard SIM coefficients; $R^2_{i,\real}$ is the real-data
coefficient of determination, used by the branch-selection step in
Section~\ref{sec:method-r2}.

\subsection{Preservation criteria}
\label{sec:method-criteria}

We want the synthetic composition
$g_i = \alpha_i + \beta_i^{\eff}\, g_m + \eps_i$, with effective slope
$\beta_i^{\eff}$ (equal to $\beta_i$ on the unclipped branch and the
attenuated value of Eq.~\eqref{eq:clip} otherwise), to satisfy three
properties in the large-$T$ limit. The first is \emph{SIM recovery}:
ordinary least squares of $g_i$ on $g_m$ returns the calibrated
intercept and slope, $\hat\alpha_i = \alpha_i$ and
$\hat\beta_i = \beta_i^{\eff}$. On the unclipped branch,
$\beta_i^{\eff} = \beta_i$ and the OLS output reproduces the
calibration exactly; on the clipped branch, $\beta_i^{\eff}$ is the
attenuated value defined by Eq.~\eqref{eq:clip}, and the attenuation
is persisted alongside the output so it is auditable. The second is
\emph{marginal fidelity}: each asset is assigned one of two per-asset
targets based on its calibrated $R^2$, not both at once. Assets whose
generator marginal (heavy tails, regime structure, jumps) is the
quantity of interest receive the variance-preserving target
$\Var[g_i] = \sigma_{\gen,i}^2$; assets that track the market closely
receive the $R^2$-preserving target $R^2 = R^2_{i,\real}$, measured by
the same OLS regression as above. The branch chosen for each asset is
recorded as a construction flag so downstream consumers can audit
which target was pursued. The third is \emph{cross-sectional
dependence}: any joint dependence structure applied to the residuals
$\{\eps_i\}_{i=1}^N$ passes through unchanged to the composed series
$\{g_i\}_{i=1}^N$, automatic given the linearity of
Eq.~\eqref{eq:sim-composition} in $\eps_i$. The joint structure
itself (empirical copula, factor copula, t-copula) is supplied as a
separate design choice, and we report cross-sectional fidelity for
the empirical-copula choice in Section~\ref{sec:results}.

The naive composition $g_i = \alpha_i + \beta_i\, g_m + \teps_i$
satisfies SIM recovery and cross-sectional dependence but violates
marginal fidelity: its marginal variance is
$\sigma_{\gen,i}^2 + \beta_i^2 \sigma_m^2$, an inflation that grows
quadratically in $\beta_i$. The rest of this section fixes this inflation with a closed-form
scalar correction, handles the degenerate high-$\beta$ case with a
sign-preserving clip, treats index-tracking assets through an
$R^2$-preserving branch, and ties the pieces together in an algorithm
that also states the resulting properties of the composed series.

\subsection{Variance correction}
\label{sec:method-variance}

Form the candidate residual $\eps_i = s_i\,\teps_i$ for some scalar
$s_i \in (0, 1]$, and require that the composed series have the same
variance as the generator,
\begin{equation}
  \Var[\beta_i\, g_m + \eps_i] \;=\; \sigma_{\gen,i}^2.
  \label{eq:variance-target}
\end{equation}
Under the independence assumption of
Section~\ref{sec:method-setup}, the cross term in
\eqref{eq:variance-target} vanishes and
\begin{equation}
  \beta_i^2\, \sigma_m^2 \;+\; s_i^2\, \sigma_{\gen,i}^2 \;=\; \sigma_{\gen,i}^2,
\end{equation}
giving
\begin{equation}
  \boxed{\;s_i^2 \;=\; 1 \;-\; \frac{\beta_i^2\, \sigma_m^2}{\sigma_{\gen,i}^2}\;}
  \label{eq:scale}
\end{equation}
Define the \emph{market loading ratio}
\begin{equation}
  \rho_i \;\equiv\; \frac{\beta_i^2\, \sigma_m^2}{\sigma_{\gen,i}^2},
  \label{eq:rho}
\end{equation}
the share of the generator variance consumed by the market component. Then
$s_i^2 = 1 - \rho_i$, valid whenever $\rho_i < 1$. The naive composition
$g_i = \alpha_i + \beta_i g_m + \teps_i$ corresponds to $s_i = 1$; it
inflates $\Var[g_i]$ by exactly $\beta_i^2 \sigma_m^2$, so the relative
distortion grows quadratically in $\beta_i$.

\subsection{Beta clipping with an idiosyncratic floor}
\label{sec:method-clipping}

Equation \eqref{eq:scale} breaks down when $\rho_i \ge 1$: the market loading
alone would account for more than the entire generator variance, leaving no
room for idiosyncratic noise. This case arises for high-$\beta_i$ assets
paired with low-volatility generators, or whenever the synthetic market
$g_m$ is more volatile than the market on which $\beta_i$ was calibrated.

To handle it, impose a floor $f \in (0, 1)$ on the idiosyncratic share by
requiring $s_i^2 \ge f$, and clip $\beta_i$ downward whenever the floor
would be violated. Writing the clip threshold as $\rho_i^{\max} = 1 - f$,
\begin{equation}
  \beta_i^{\eff} \;=\;
  \begin{cases}
    \beta_i & \text{if } \rho_i \le \rho_i^{\max} \\[4pt]
    \sign(\beta_i)\,\sqrt{(1 - f)\,\dfrac{\sigma_{\gen,i}^2}{\sigma_m^2}} & \text{if } \rho_i > \rho_i^{\max}
  \end{cases}
  \qquad
  s_i^2 \;=\;
  \begin{cases}
    1 - \rho_i & \text{if } \rho_i \le \rho_i^{\max} \\[4pt]
    f          & \text{if } \rho_i > \rho_i^{\max}
  \end{cases}
  \label{eq:clip}
\end{equation}
Sign preservation in \eqref{eq:clip} keeps defensive (negative-$\beta$)
assets defensive even when clipped. The floor $f$ is a tunable knob; in the
experiments below we use $f = 0.10$, which guarantees that at least
$10\%$ of each asset's variance comes from idiosyncratic noise. When
$\beta_i$ is clipped we log a warning and persist both $\beta_i$ and
$\beta_i^{\eff}$ so downstream consumers can audit the attenuation.

\subsection{$R^2$-preserving branch for tracker assets}
\label{sec:method-r2}

The variance-preserving construction targets $\Var[g_i] = \sigma_{\gen,i}^2$.
That is the right target when the generator marginal carries information one
wants to keep. For some assets the generator marginal is not really the
quantity of interest; the SIM relationship to the market is. Index ETFs are
the canonical example. SPY against itself has $\beta = 1$, $R^2 = 1$,
$\sigma_\eps = 0$ exactly. QQQ against SPY has $R^2 \approx 0.84$, SPYG
against SPY has $R^2 \approx 0.86$. For these tickers the practitioner cares that $\hat\beta$ recovers
to its calibrated value and that $R^2$ recovers to its calibrated
value, neither inflated to one by dropping $\eps$ nor implied by
unrelated marginal-variance bookkeeping.

We branch on the real-data calibration $R^2_{i,\real}$. Pick a threshold
$R^2_{\mathrm{preserve}}$ that separates index-tracking assets from
genuine single stocks; in our universe the cumulative distribution sorts
itself cleanly, with index ETFs above $0.80$ and the most market-like
single names below $0.65$, so $R^2_{\mathrm{preserve}} = 0.80$ captures the
trackers without sweeping in any genuinely idiosyncratic name. For any
asset with $R^2_{i,\real} \ge R^2_{\mathrm{preserve}}$, the population
$R^2$ identity for \eqref{eq:sim-composition},
\begin{equation}
  R^2 \;=\; \frac{\beta_i^2\, \sigma_m^2}{\beta_i^2\, \sigma_m^2 + \sigma_{\eps,i}^2},
\end{equation}
inverts to give the residual variance that hits the calibrated $R^2$:
\begin{equation}
  \boxed{\;\sigma_{\eps,\mathrm{target}}^2 \;=\; \beta_i^2\, \sigma_m^2 \cdot
    \frac{1 - R^2_{i,\real}}{R^2_{i,\real}}\;}
  \label{eq:r2-target}
\end{equation}
Draw $\teps_i$ from the ticker's generator and rescale,
\begin{equation}
  \eps_i \;=\; \sqrt{\frac{\sigma_{\eps,\mathrm{target}}^2}{\sigma_{\gen,i}^2}}\,\teps_i,
  \label{eq:r2-rescale}
\end{equation}
then copula-reorder and compose as in \eqref{eq:sim-composition}. The
construction recovers both calibrated quantities by design:
$\hat\beta \to \beta_i$ in the large-$T$ limit because $\eps_i$ is
uncorrelated with $g_m$ after the rank reorder, and $R^2 \to R^2_{i,\real}$
by construction.

When $R^2_{i,\real} = 1$ exactly (SPY against itself), \eqref{eq:r2-target}
gives $\sigma_{\eps,\mathrm{target}}^2 = 0$, so $\eps_i \equiv 0$ and
$g_i = \alpha_i + \beta_i g_m$ deterministically. No special-casing is
needed: the SPY identity falls out of the same expression with
$\beta = 1$, $\alpha = 0$, $R^2 = 1$. If the rescale factor in
\eqref{eq:r2-rescale} exceeds one, the generator's tails are being stretched
to fit a target larger than the original draw; this happens when the
synthetic market $g_m$ has lower variance than the calibration market.
We detect and flag this case so the loss of marginal fidelity is auditable.

\subsection{Branch selection, algorithm, and properties}
\label{sec:method-algorithm}

The pieces combine into a single procedure (Algorithm~\ref{alg:hybrid}).
For each asset the calibrated $R^2_{i,\real}$ selects a branch: assets with
$R^2_{i,\real} \ge R^2_{\mathrm{preserve}}$ go through the
$R^2$-preserving construction of Section~\ref{sec:method-r2}, and the
rest go through the variance-preserving construction of
Section~\ref{sec:method-variance} with clipping
(Section~\ref{sec:method-clipping}) engaged when the market loading
ratio $\rho_i$ would otherwise exceed $1 - f$. Each asset is tagged
with a construction flag $\texttt{r2-preserve}$, $\texttt{hybrid}$,
or $\texttt{hybrid-clipped}$ that is persisted alongside the output,
so downstream consumers can audit which path each asset took. The optional rank reorder in line~14 accepts any
joint dependence structure (empirical copula, t-copula, factor
copula); reordering preserves the marginal variance of $\eps_i$, so
the variance algebra of Section~\ref{sec:method-variance} is
unaffected and whatever cross-sectional structure is applied to
$\eps_i$ passes through to $g_i$ unchanged.

\begin{algorithm}[t]
\caption{Hybrid SIM composition with variance correction.}
\label{alg:hybrid}
\begin{algorithmic}[1]
  \REQUIRE Market path $g_m \in \R^T$ with sample variance $\sigma_m^2$;
    per-asset generator draws $\teps_i \in \R^T$ for $i = 1,\dots,N$;
    calibrated SIM parameters $(\alpha_i, \beta_i, R^2_{i,\real})$;
    floor $f \in (0,1)$; threshold $R^2_{\mathrm{preserve}} \in (0,1)$.
  \ENSURE Composed paths $g_i \in \R^T$ and per-asset construction flags.
  \FOR{each asset $i = 1,\dots,N$}
    \STATE Compute generator variance $\sigma_{\gen,i}^2 \gets \Var[\teps_i]$.
    \IF{$R^2_{i,\real} \ge R^2_{\mathrm{preserve}}$}
      \STATE $\sigma_{\eps,\mathrm{target}}^2 \gets \beta_i^2 \sigma_m^2 (1 - R^2_{i,\real})/R^2_{i,\real}$
        \hfill \COMMENT{Eq.~\eqref{eq:r2-target}}
      \STATE $\eps_i \gets \sqrt{\sigma_{\eps,\mathrm{target}}^2 / \sigma_{\gen,i}^2}\,\teps_i$
      \STATE $\beta_i^{\eff} \gets \beta_i$;\quad flag $\gets$ \texttt{r2-preserve}
    \ELSE
      \STATE $\rho_i \gets \beta_i^2 \sigma_m^2 / \sigma_{\gen,i}^2$
        \hfill \COMMENT{Eq.~\eqref{eq:rho}}
      \IF{$\rho_i \le 1 - f$}
        \STATE $s_i^2 \gets 1 - \rho_i$;\quad $\beta_i^{\eff} \gets \beta_i$;\quad flag $\gets$ \texttt{hybrid}
      \ELSE
        \STATE $s_i^2 \gets f$;\quad
          $\beta_i^{\eff} \gets \sign(\beta_i)\sqrt{(1-f)\,\sigma_{\gen,i}^2/\sigma_m^2}$;\quad
          flag $\gets$ \texttt{hybrid-clipped}
      \ENDIF
      \STATE $\eps_i \gets s_i\,\teps_i$
    \ENDIF
    \STATE \textit{(Optional)} apply cross-sectional rank reorder to $\eps_i$.
    \STATE $g_i \gets \alpha_i + \beta_i^{\eff}\, g_m + \eps_i$
      \hfill \COMMENT{Eq.~\eqref{eq:sim-composition}}
  \ENDFOR
\end{algorithmic}
\end{algorithm}

By construction the composed series $g_i$ satisfies the preservation
criteria of Section~\ref{sec:method-criteria} in the large-$T$ limit.
OLS of $g_i$ on $g_m$ recovers $\hat\alpha_i \to \alpha_i$ and
$\hat\beta_i \to \beta_i^{\eff}$, with $\beta_i^{\eff} = \beta_i$
unclipped and $\beta_i^{\eff}$ persisted alongside $\beta_i$ when the
clip triggers. Marginal variance matches $\sigma_{\gen,i}^2$ on the
variance-preserving branch, either exactly,
$\Var[g_i] = \beta_i^2 \sigma_m^2 + s_i^2 \sigma_{\gen,i}^2
  = \sigma_{\gen,i}^2$, or by the choice of $\beta_i^{\eff}$ in
\eqref{eq:clip} under clipping; on the $R^2$-preserving branch the
marginal variance instead matches
$\beta_i^2 \sigma_m^2 / R^2_{i,\real}$, the value implied by the
calibrated $(\beta_i, R^2_{i,\real})$ rather than by the generator
marginal. Tail behavior on the variance-preserving branch follows
from the floor: because at least $f\,\sigma_{\gen,i}^2$ of the
variance is reserved for $\eps_i$, the tails inherited from $\teps_i$
dominate whenever $(\beta_i^{\eff})^2 \sigma_m^2$ is small relative
to those tails, while as $\beta_i$ grows the marginal of $g_i$
tightens toward Gaussian along the market-driven axis, mildly for
low-$\beta$ assets and noticeably at high $\beta$. The construction
injects \emph{contemporaneous} market dependence only; lagged
dependence (lead-lag effects, market-driven volatility clustering)
requires a richer construction, which we return to in
Section~\ref{sec:discussion}. The pipeline is modest in compute:
fitting the $423$ per-asset JumpHMM marginals runs once offline
against the $2014$--$2024$ price histories; the scoring loop across
$423 \times 100 = 42{,}300$ per-ticker replications per composer
runs in minutes on a commodity laptop, and the full six-composer
sweep reported in Section~\ref{sec:results} completes within an
hour on the same hardware.
